% Options for packages loaded elsewhere
\PassOptionsToPackage{unicode}{hyperref}
\PassOptionsToPackage{hyphens}{url}
\PassOptionsToPackage{dvipsnames,svgnames,x11names}{xcolor}
\documentclass[
  10pt,
]{article}
\usepackage{xcolor}
\usepackage[margin=1cm,top=1cm,bottom=1cm,left=1cm,right=1cm,includeheadfoot]{geometry}
\usepackage{amsmath,amssymb}
\setcounter{secnumdepth}{3}
\usepackage{iftex}
\ifPDFTeX
  \usepackage[T1]{fontenc}
  \usepackage[utf8]{inputenc}
  \usepackage{textcomp} % provide euro and other symbols
\else % if luatex or xetex
  \usepackage{unicode-math} % this also loads fontspec
  \defaultfontfeatures{Scale=MatchLowercase}
  \defaultfontfeatures[\rmfamily]{Ligatures=TeX,Scale=1}
\fi
\usepackage{lmodern}
\ifPDFTeX\else
  % xetex/luatex font selection
  \setmainfont[]{DejaVu Serif}
  \setmonofont[]{DejaVu Sans Mono}
\fi
% Use upquote if available, for straight quotes in verbatim environments
\IfFileExists{upquote.sty}{\usepackage{upquote}}{}
\IfFileExists{microtype.sty}{% use microtype if available
  \usepackage[]{microtype}
  \UseMicrotypeSet[protrusion]{basicmath} % disable protrusion for tt fonts
}{}
\usepackage{setspace}
\makeatletter
\@ifundefined{KOMAClassName}{% if non-KOMA class
  \IfFileExists{parskip.sty}{%
    \usepackage{parskip}
  }{% else
    \setlength{\parindent}{0pt}
    \setlength{\parskip}{6pt plus 2pt minus 1pt}}
}{% if KOMA class
  \KOMAoptions{parskip=half}}
\makeatother
\usepackage{listings}
\newcommand{\passthrough}[1]{#1}
\lstset{defaultdialect=[5.3]Lua}
\lstset{defaultdialect=[x86masm]Assembler}
\setlength{\emergencystretch}{3em} % prevent overfull lines
\providecommand{\tightlist}{%
  \setlength{\itemsep}{0pt}\setlength{\parskip}{0pt}}
% Enable graphics inclusion and ensure figure numbering works
\usepackage{graphicx}
\renewcommand{\figurename}{Figure}

% Configure fonts for Unicode support with fallbacks
\usepackage{newunicodechar}
\newunicodechar{⁴}{\textsuperscript{4}}
\newunicodechar{₄}{\textsubscript{4}}

% Math notation macros
\newcommand{\norm}[1]{\lVert #1\rVert}
\newcommand{\abs}[1]{\lvert #1\rvert}

% Enhanced code block styling for better contrast and readability
\usepackage{fancyvrb}
\usepackage{xcolor}
\usepackage{listings}

% Define custom colors for code blocks
\definecolor{codebg}{RGB}{245, 245, 245}      % Light gray background
\definecolor{codeborder}{RGB}{200, 200, 200}  % Medium gray border
\definecolor{codefg}{RGB}{50, 50, 50}         % Dark gray text

% Configure Verbatim environment for inline code
\DefineVerbatimEnvironment{Verbatim}{Verbatim}{%
    fontsize=\small,
    frame=single,
    framerule=0.5pt,
    framesep=3pt,
    rulecolor=\color{codeborder},
    bgcolor=\color{codebg},
    fgcolor=\color{codefg}
}

% Configure code block styling
\DefineVerbatimEnvironment{Highlighting}{Verbatim}{%
    fontsize=\footnotesize,
    frame=single,
    framerule=0.5pt,
    framesep=5pt,
    rulecolor=\color{codeborder},
    bgcolor=\color{codebg},
    fgcolor=\color{codefg}
}

% Style inline code with \texttt
\renewcommand{\texttt}[1]{%
    \colorbox{codebg}{\color{codefg}\ttfamily #1}%
}

% Configure listings package for code blocks
\lstset{
    backgroundcolor=\color{codebg},
    basicstyle=\footnotesize\ttfamily\color{codefg},
    breakatwhitespace=false,
    breaklines=true,
    captionpos=b,
    commentstyle=\color{codefg},
    deletekeywords={...},
    escapeinside={\%*}{*)},
    extendedchars=true,
    frame=single,
    framerule=0.5pt,
    framesep=5pt,
    keepspaces=true,
    keywordstyle=\color{codefg},
    language=Python,
    morekeywords={*,...},
    numbers=left,
    numbersep=5pt,
    numberstyle=\tiny\color{codefg},
    rulecolor=\color{codeborder},
    showspaces=false,
    showstringspaces=false,
    showtabs=false,
    stepnumber=1,
    stringstyle=\color{codefg},
    tabsize=2,
    title=\lstname
}

% Override any Pandoc default lstset configurations
\AtBeginDocument{
    \lstset{
        backgroundcolor=\color{codebg},
        basicstyle=\footnotesize\ttfamily\color{codefg},
        frame=single,
        framerule=0.5pt,
        framesep=5pt,
        rulecolor=\color{codeborder},
        numbers=left,
        numbersep=5pt,
        numberstyle=\tiny\color{codefg}
    }
}

% Configure hyperref colors consistently
\AtBeginDocument{
% Override pandoc's hidelinks setting with consistent options
\hypersetup{
    colorlinks=true,
    allcolors=red,
    linkcolor=red,
    urlcolor=red,
    citecolor=red,
    filecolor=red,
    menucolor=red,
    linktoc=all
}
}

% Simple page break support for document structure
\usepackage{bookmark}
\IfFileExists{xurl.sty}{\usepackage{xurl}}{} % add URL line breaks if available
\urlstyle{same}
% fallback for those not using the hyperref driver hyperxmp:
\makeatletter
\@ifundefined{xmpquote}{\newcommand{\xmpquote}[1]{#1}}{}
\makeatother
\hypersetup{
  colorlinks=true,
  linkcolor={red},
  filecolor={red},
  citecolor={red},
  urlcolor={red},
  pdfcreator={LaTeX via pandoc}}

\title{Equations and Math Supplement}
\author{Daniel Ari Friedman\\ ORCID: 0000-0001-6232-9096\\ Email: daniel@activeinference.institute\\ DOI: 10.5281/zenodo.16887791}
\date{October 07, 2026}

\begin{document}
\maketitle

{
\hypersetup{linkcolor=black}
\setcounter{tocdepth}{3}
\tableofcontents
}
\setstretch{1.0}
\section{Equations and Math Supplement
(Appendix)}\label{equations-and-math-supplement-appendix}

Conventions for this appendix: vertices are \(P_0,\ldots,P_3\) with
Quadray components \((a_i, b_i, c_i, d_i)\); volumes are \(V_{xyz}\)
(Euclidean, cubic length units) and \(V_{ivm}\) (synergetics/IVM
tetravolume units, unit regular tetrahedron \(V_{ivm} = 1\)), related by
\(V_{ivm} = S3\,V_{xyz}\) with \(S3=\sqrt{9/8}\) (Section 3); \(M\)
denotes the Quadray-to-XYZ embedding matrix (Sections 3 and 14),
\(\lVert\cdot\rVert_2\) the Euclidean norm, and
\(\mathrm{KL}\big[Q\,\|\,P\big]\) the Kullback--Leibler divergence of
\(Q\) relative to \(P\) (Section 10). Three symbol overloads are scoped
by context and flagged where they occur: \(c\) is a Quadray component
(and a PdF edge length) in the volume formulas but the speed of light in
the Minkowski line element; \(d\) is the fourth Quadray component in the
coordinate formulas but an edge length in the length-based formulas and
the distance function \(d(\cdot,\cdot)\) of the embedding section; and
\(P\) is a vertex label \(P_i\) in the volume sections but a probability
distribution in the free-energy sections.

\subsection{Volume of a Tetrahedron
(Lattice)}\label{volume-of-a-tetrahedron-lattice}

\begin{equation}\label{eq:lattice_det}
V_{xyz} = \tfrac{1}{6}\,\left|\det\,[\,P_1 - P_0,\; P_2 - P_0,\; P_3 - P_0\,]\right|
\end{equation}

Notes.

\begin{itemize}
\tightlist
\item
  \(P_0,\ldots,P_3\) are Cartesian (XYZ) vertex coordinates, in length
  units; each bracketed column is an edge vector, the determinant is the
  volume of the parallelepiped they span, and the \(1/6\) factor
  converts it to the tetrahedron volume \(V_{xyz}\) in cubic length
  units. This is the coordinate (difference) form of the homogeneous-row
  determinant of Eq. \eqref{eq:xyz_det}; for IVM units directly from
  Quadray coordinates, see the native formula of Eq. \eqref{eq:gdj}.
\end{itemize}

Tom Ace 5×5 tetravolume (IVM units):

\begin{equation}\label{eq:ace5x5}
V_{ivm} = \tfrac{1}{4} \left| \det \begin{pmatrix}
 a_0 & b_0 & c_0 & d_0 & 1 \\
 a_1 & b_1 & c_1 & d_1 & 1 \\
 a_2 & b_2 & c_2 & d_2 & 1 \\
 a_3 & b_3 & c_3 & d_3 & 1 \\
  1 & 1 & 1 & 1 & 0
\end{pmatrix} \right|
\end{equation}

Notes.

\begin{itemize}
\tightlist
\item
  Row \(i\) lists the four Quadray components \((a_i, b_i, c_i, d_i)\)
  of vertex \(P_i\), augmented with an affine 1; the last row
  \((1,1,1,1,0)\) encodes the projective normalization constraint, so
  the determinant is invariant to adding \((t,t,t,t)\) to every vertex
  (Section 3). Division by 4 returns the IVM tetravolume \(V_{ivm}\);
  for integer quadrays the determinant is an exact integer (Bareiss
  algorithm), computed as a
  \passthrough{\lstinline!fractions.Fraction!}. This matrix is
  identical, row for row, to the named matrix \(\mathsf{Q}\) of Eq.
  \eqref{eq:ace5x5_expanded}.
\end{itemize}

\subsection{Expanded Ace 5×5 Matrix}\label{expanded-ace-55-matrix}

The Ace 5×5 matrix of Eq. \eqref{eq:ace5x5} written out explicitly and
named \(\mathsf{Q}\):

\begin{equation}\label{eq:ace5x5_expanded}
\mathsf{Q}(P_0,P_1,P_2,P_3) = \begin{bmatrix}
 a_0 & b_0 & c_0 & d_0 & 1 \\
 a_1 & b_1 & c_1 & d_1 & 1 \\
 a_2 & b_2 & c_2 & d_2 & 1 \\
 a_3 & b_3 & c_3 & d_3 & 1 \\
1 & 1 & 1 & 1 & 0
\end{bmatrix}, \qquad V_{ivm} = \tfrac{1}{4}\,\big|\det \mathsf{Q}(P_0,\ldots,P_3)\big|
\end{equation}

Notes.

\begin{itemize}
\tightlist
\item
  \textbf{Matrix structure}: row \(i\) holds the four Quadray components
  \((a_i, b_i, c_i, d_i)\) of vertex \(P_i\), plus the affine coordinate
  1 --- literally the same rows as Eq. \eqref{eq:ace5x5}.
\item
  \textbf{Last row}: \((1,1,1,1,0)\) enforces the projective
  normalization constraint (Section 3).
\item
  \textbf{Volume computation}:
  \(V_{ivm} = \tfrac{1}{4}\,\big|\det \mathsf{Q}\big|\) in IVM units,
  exact for integer quadrays.
\item
  \textbf{Notation}: the Ace matrix is written \(\mathsf{Q}\); the
  symbol \(M\) is reserved for the Quadray-to-XYZ embedding matrix
  (Sections 3 and 14).
\end{itemize}

XYZ determinant volume and S3 conversion:

\begin{equation}\label{eq:xyz_det}
V_{xyz} = \tfrac{1}{6} \left| \det \begin{pmatrix}
 x_0 & y_0 & z_0 & 1 \\
 x_1 & y_1 & z_1 & 1 \\
 x_2 & y_2 & z_2 & 1 \\
 x_3 & y_3 & z_3 & 1 \\
\end{pmatrix} \right|, \qquad V_{ivm} = S3\, V_{xyz},\quad S3=\sqrt{\tfrac{9}{8}}
\end{equation}

Notes.

\begin{itemize}
\tightlist
\item
  Homogeneous-row determinant in Cartesian coordinates: each row is
  \((x_i, y_i, z_i, 1)\) for vertex \(P_i\), with coordinates in length
  units; the absolute determinant divided by 6 is the Euclidean volume
  \(V_{xyz}\) in cubic length units. This is the homogeneous form of Eq.
  \eqref{eq:lattice_det}; conversion to IVM units uses
  \(V_{ivm} = S3\,V_{xyz}\) with \(S3=\sqrt{9/8}\) as used throughout.
\end{itemize}

\subsection{Cayley-Menger Determinant
(Coxeter.4D)}\label{cayley-menger-determinant-coxeter.4d}

For tetrahedron volume from edge lengths (Coxeter.4D approach):

\begin{equation}\label{eq:cayley_menger}
288\,V_{xyz}^2 = \det\begin{pmatrix}
  0 & 1 & 1 & 1 & 1 \\
  1 & 0 & d_{01}^2 & d_{02}^2 & d_{03}^2 \\
  1 & d_{10}^2 & 0 & d_{12}^2 & d_{13}^2 \\
  1 & d_{20}^2 & d_{21}^2 & 0 & d_{23}^2 \\
  1 & d_{30}^2 & d_{31}^2 & d_{32}^2 & 0
\end{pmatrix}
\end{equation}

Notes.

\begin{itemize}
\tightlist
\item
  \textbf{Pairwise distances}: \(d_{ij}\) is the Euclidean distance
  between vertices \(P_i\) and \(P_j\) in length units; the matrix
  stores the squared distances \(d_{ij}^2\), so the formula is
  length-only --- no coordinates enter (Coxeter.4D).
\item
  \textbf{Length-only formulation}: Cayley--Menger provides a
  length-only formula for simplex volumes, here specialized to
  tetrahedra.
\item
  \textbf{Conversion to IVM}: \(V_{xyz}\) is the Euclidean volume in
  cubic length units; use \(V_{ivm} = S3\,V_{xyz}\) with
  \(S3=\sqrt{9/8}\) (Eq. \eqref{eq:xyz_det}); the PdF formula of Eq.
  \eqref{eq:pdf} consumes the same lengths in closed algebraic form.
\end{itemize}

\subsection{Piero della Francesca Formula
(PdF)}\label{piero-della-francesca-formula-pdf}

For tetrahedron volume from edge lengths meeting at a vertex:

\begin{equation}\label{eq:pdf}
144\,V_{xyz}^2 = 4 a^2 b^2 c^2 - a^2\,(b^2 + c^2 - f^2)^2 - b^2\,(c^2 + a^2 - e^2)^2 - c^2\,(a^2 + b^2 - d^2)^2 + (b^2 + c^2 - f^2)(c^2 + a^2 - e^2)(a^2 + b^2 - d^2)
\end{equation}

Notes.

\begin{itemize}
\tightlist
\item
  \textbf{Edge lengths}: \(a, b, c\) are the three edges meeting at the
  apex vertex \(P_0\) --- \(a = d_{01}\), \(b = d_{02}\), \(c = d_{03}\)
  --- and \(d, e, f\) are the respective opposite edges, \(d = d_{23}\),
  \(e = d_{13}\), \(f = d_{12}\), in length units (the same \(d_{ij}\)
  as Eq. \eqref{eq:cayley_menger}).
\item
  \textbf{Conversion to IVM}: \(V_{xyz}\) is the Euclidean volume in
  cubic length units (the prefactor 144 is the standard Heron-like PdF
  normalization); use \(V_{ivm} = S3\,V_{xyz}\) with \(S3=\sqrt{9/8}\)
  (Eq. \eqref{eq:xyz_det}).
\end{itemize}

\subsection{Gerald de Jong Formula
(GdJ)}\label{gerald-de-jong-formula-gdj}

Native Quadray formula for tetrahedron volume:

\begin{equation}\label{eq:gdj}
V_{ivm} = \tfrac{1}{4}\,\left|\det\big[\, \pi(P_1) - \pi(P_0),\; \pi(P_2) - \pi(P_0),\; \pi(P_3) - \pi(P_0) \,\big]\right|, \qquad \pi(P_i) = (\,a_i - d_i,\; b_i - d_i,\; c_i - d_i\,)
\end{equation}

Notes.

\begin{itemize}
\tightlist
\item
  \textbf{Projection}: \(\pi\) maps a Quadray vertex \(P_i\) with
  components \((a_i, b_i, c_i, d_i)\) to \(\mathbb{R}^3\) by subtracting
  the fourth component from the first three; each bracketed column is a
  projected edge vector. Because \(\pi\) kills the direction
  \((1,1,1,1)\), it is invariant under \(q \to q + t\,(1,1,1,1)\) ---
  the projective fiber of Eq. \eqref{eq:conv-fiber} (Section 14) --- so
  the determinant is unchanged by projective normalization of the
  vertices. When all four vertices share the same fourth component (in
  particular \(d_i = 0\)), \(\pi(P_i)\) reduces to \((a_i, b_i, c_i)\).
\item
  \textbf{Native IVM}: no S3 conversion; the factor \(1/4\) returns IVM
  tetravolume directly --- for the unit tetrahedron (origin plus three
  IVM neighbor moves) the determinant is exactly 4, giving
  \(V_{ivm} = 1\) (Section 3).
\item
  \textbf{Exact arithmetic}: integer Quadray coordinates give an integer
  determinant, evaluated exactly by the Bareiss algorithm as a
  \passthrough{\lstinline!fractions.Fraction!}; the code implements Eq.
  \eqref{eq:gdj} as \passthrough{\lstinline!integer\_tetra\_volume!},
  and the magnitude agrees exactly with the Ace determinant of Eq.
  \eqref{eq:ace5x5} on every integer-quadray input.
\end{itemize}

See code:
\href{03_quadray_methods.md\#code:tetra_volume_cayley_menger}{\passthrough{\lstinline!tetra\_volume\_cayley\_menger!}}.
For tetrahedron volume background, see
\href{https://en.wikipedia.org/wiki/Tetrahedron\#Volume}{Tetrahedron --
volume}. Exact integer determinants in code use the
\href{https://en.wikipedia.org/wiki/Bareiss_algorithm}{Bareiss
algorithm}. External validation: these formulas align with
implementations in the 4dsolutions ecosystem. See the
\href{07_resources.md}{Resources} section for comprehensive details.

\subsection{Fisher Information Matrix (FIM)}\label{eq:fim}

Background:
\href{https://en.wikipedia.org/wiki/Fisher_information}{Fisher
information}.

\begin{equation}\label{eq:fim}
F_{i,j} = \mathbb{E}\left[ \frac{\partial \, \log p(x;\theta)}{\partial \theta_i}\, \frac{\partial \, \log p(x;\theta)}{\partial \theta_j} \right]
\end{equation}

Notes.

\begin{itemize}
\tightlist
\item
  \textbf{Symbols}: \(x\) is an observation,
  \(\theta = (\theta_1, \ldots, \theta_m)\) the parameter vector, and
  \(p(x;\theta)\) the likelihood; the expectation is over
  \(x \sim p(\cdot;\theta)\), so \(F\) is a function of \(\theta\). Each
  score \(\partial \log p(x;\theta)/\partial \theta_i\) carries inverse
  parameter units, so \(F_{i,j}\) carries inverse squared parameter
  units per observation.
\item
  \(F\) is symmetric positive semi-definite: large eigenvalues mark
  sensitive (high-curvature) directions, small eigenvalues sloppy ones
  (Section 4 and the Discussion). The empirical estimate of Eq.
  \eqref{eq:fim_empirical} is the per-observation counterpart.
\item
  See code:
  \href{03_quadray_methods.md\#code:fisher_information_matrix}{\passthrough{\lstinline!fisher\_information\_matrix!}}
  in \passthrough{\lstinline!src/quadmath/inference/information.py!} ---
  empirical outer-product estimator. Figure: the empirical estimate is
  shown in the FIM heatmap figure of Section 4.
\end{itemize}

\subsection{Empirical Fisher Information
Matrix}\label{empirical-fisher-information-matrix}

For empirical estimation from data, the Fisher Information Matrix is
computed as:

\begin{equation}\label{eq:fim_empirical}
F_{i,j} = \frac{1}{N} \sum_{n=1}^{N} \frac{\partial \, \log p(x_n;\theta)}{\partial \theta_i}\, \frac{\partial \, \log p(x_n;\theta)}{\partial \theta_j}
\end{equation}

Notes.

\begin{itemize}
\tightlist
\item
  \textbf{Symbols}: \(x_1, \ldots, x_N\) are \(N\) i.i.d. observations
  and \(g_n = \nabla_\theta \log p(x_n;\theta)\) the per-sample score
  vector (same units as Eq. \eqref{eq:fim}); the estimator is
  \(F = \tfrac{1}{N} \sum_n g_n g_n^{\mathsf{T}}\), with a
  \passthrough{\lstinline!normalize!} flag controlling the \(1/N\)
  division.
\item
  Converges to Eq. \eqref{eq:fim} as \(N \to \infty\); used by
  natural-gradient descent (Eq. \eqref{eq:natural_gradient}) and the
  information-geometry applications of Section 4.
\end{itemize}

\subsection{Natural Gradient}\label{eq:natgrad}

Background:
\href{https://en.wikipedia.org/wiki/Natural_gradient}{Natural gradient}
(Amari).

\begin{equation}\label{eq:natural_gradient}
\theta \leftarrow \theta - \eta\, F(\theta)^{-1}\, \nabla_{\theta} L(\theta)
\end{equation}

Explanation.

\begin{itemize}
\tightlist
\item
  \textbf{Symbols}: \(\theta\) the parameter vector, \(L(\theta)\) a
  differentiable loss, \(\nabla_\theta L\) its gradient, \(F(\theta)\)
  the Fisher matrix of Eq. \eqref{eq:fim} --- or its empirical estimate,
  Eq. \eqref{eq:fim_empirical} --- evaluated at \(\theta\), and
  \(\eta > 0\) the step size (learning rate).
\item
  \textbf{Update}: right-preconditioning by the inverse Fisher metric
  gives the steepest-descent direction under the Fisher metric (Amari),
  invariant to smooth invertible reparameterizations of \(\theta\).
\item
  \textbf{Damping}: the implementation
  \passthrough{\lstinline!natural\_gradient\_step!} solves the damped
  system \((F + \lambda I)\,\delta = \nabla_\theta L\) and applies
  \(\theta \leftarrow \theta - \eta\,\delta\), with a small Tikhonov
  ridge \(\lambda\) (default \(10^{-9}\)) keeping the solve stable when
  \(F\) is near-singular.
\end{itemize}

See code:
\href{03_quadray_methods.md\#code:natural_gradient_step}{\passthrough{\lstinline!natural\_gradient\_step!}}
in \passthrough{\lstinline!src/quadmath/inference/information.py!} ---
damped inverse-Fisher step.

\subsection{Free Energy (Active Inference)}\label{eq:free_energy}

\begin{equation}\label{eq:free_energy}
\mathcal{F} = -\log P(o\mid s) + \mathrm{KL}\big[ Q(s)\;\|\; P(s) \big]
\end{equation}

Explanation.

\begin{itemize}
\tightlist
\item
  \textbf{Symbols}: \(o\) the observation, \(s\) the latent state,
  \(Q(s)\) the approximate posterior (recognition distribution), and
  \(P\) the generative model, so \(P(o \mid s)\) is the likelihood and
  \(P(s)\) the prior; \(\mathcal{F}\) is the variational free energy in
  nats (natural logarithm, Section 10). \textbf{Notation flag}: capital
  \(P\) here is a probability distribution, unrelated to the vertex
  labels \(P_i\) of the volume sections.
\item
  \textbf{Partition}: minimizing \(\mathcal{F}\) over \(Q\) trades the
  expected negative log-likelihood \(\mathbb{E}_Q[-\log P(o \mid s)]\)
  (the displayed \(-\log P(o\mid s)\) is read under this
  \(Q\)-expectation, as implemented) against the KL divergence
  \(\mathrm{KL}\big[Q\,\|\,P\big]\) that ties the posterior to the
  prior; see
  \href{https://en.wikipedia.org/wiki/Free_energy_principle}{Free energy
  principle}.
\end{itemize}

See code:
\href{03_quadray_methods.md\#code:free_energy}{\passthrough{\lstinline!free\_energy!}}
in \passthrough{\lstinline!src/quadmath/inference/information.py!} ---
discrete-state variational free energy (inputs are unnormalized
distributions, normalized internally).

\textbf{Note}: The main figures demonstrating natural gradient
trajectories and free energy landscapes are shown in
\href{04_optimization_in_4d.md}{Section 4: Optimization in 4D}. The
appendix focuses on unique figures specific to mathematical formulations
and validation.

\subsection{Expected Free Energy (Active
Inference)}\label{eq:expected_free_energy}

Background:
\href{https://direct.mit.edu/books/oa-monograph/5299/Active-InferenceThe-Free-Energy-Principle-in-Mind}{Active
Inference (Parr, Pezzulo \& Friston, MIT Press, 2022)}.

\begin{equation}\label{eq:expected_free_energy}
G = \mathrm{KL}\big[ Q(s)\;\|\;P(s) \big] \;-\; \mathbb{E}_{q}\big[\log P(o\mid s)\big] \;-\; \log P(o)
\end{equation}

Explanation.

\begin{itemize}
\tightlist
\item
  \textbf{Symbols}: \(Q(s)\) the approximate posterior and \(P(s)\) the
  prior over states (as in Eq. \eqref{eq:free_energy}), \(P(o \mid s)\)
  the likelihood, \(P(o)\) the prior preference over outcomes in nats
  (uniform when omitted), \(H\big[Q\big] = -\sum_s Q(s) \log Q(s)\) the
  Shannon entropy in nats, and \(\mathbb{E}_q\) expectation under \(Q\);
  \(G\) is the expected free energy minimized during action selection.
\item
  \textbf{Epistemic term}: the KL divergence between variational
  posterior and prior over states.
\item
  \textbf{Entropy}: the posterior entropy is already inside the KL term,
  since
  \(\mathrm{KL}[Q\|P] = \mathbb{E}_q[\log Q(s)] - \mathbb{E}_q[\log P(s)] = -H[Q(s)] - \mathbb{E}_q[\log P(s)]\);
  subtracting \(H\) again would count it twice.
\item
  \textbf{Ambiguity}: the negative expected log-likelihood of outcomes
  penalizes noisy observations.
\item
  \textbf{Pragmatic term}: prior preferences \(P(o)\) enter negatively
  (\(-\log P(o)\)), so preferred outcomes lower \(G\); agents minimize
  \(G\) during action selection.
\end{itemize}

See code:
\href{03_quadray_methods.md\#code:expected_free_energy}{\passthrough{\lstinline!expected\_free\_energy!}}
in \passthrough{\lstinline!src/quadmath/inference/information.py!} ---
all four terms of Eq. \eqref{eq:expected_free_energy} as implemented.

\subsection{Quadray Normalization
(Fuller.4D)}\label{quadray-normalization-fuller.4d}

Given a Quadray \(q = (a, b, c, d)\), choose \(k = \min(a, b, c, d)\)
and set \(q' = q - (k, k, k, k)\), enforcing non-negative entries with
at least one zero. The shift is the projective normalization of Section
3 --- \(q\) and \(q'\) represent the same direction, the fiber
formalized in Eq. \eqref{eq:conv-fiber} (Section 14). Here \(k\) is the
normalization offset of the glossary (Section 10); it is unrelated to
the shell-frequency \(k\) of the lattice sections (Sections 11 and 13).

\subsection{Distance (Embedding Sketch; Coxeter.4D
slice)}\label{distance-embedding-sketch-coxeter.4d-slice}

Choose a linear map \(M\) from Quadray space to \(\mathbb{R}^3\) (or
\(\mathbb{R}^4\)) consistent with the tetrahedral axes --- the embedding
matrix of Sections 3 and 14 (canonical Urner family, Eq.
\eqref{eq:conv-embedding}). Then for Quadray points \(p, q\), the
distance is \(d(p, q) = \lVert M\,(p - q) \rVert_2\) in length units;
under the integer Urner family the squared distance is the exact integer
identity of Eq. \eqref{eq:conv-distance}, which the lattice-search layer
ranks with (Section 13).

\subsection{Minkowski Line Element (Einstein.4D
analogy)}\label{minkowski-line-element-einstein.4d-analogy}

\begin{equation}\label{eq:minkowski_line_element}
ds^2 = -c^2\,dt^2 + dx^2 + dy^2 + dz^2
\end{equation}

Background:
\href{https://en.wikipedia.org/wiki/Minkowski_space}{Minkowski space}.

Signature convention \((-,+,+,+)\) (mostly-plus, Section 2): \(c\) is
the speed of light, \((t, x, y, z)\) are spacetime coordinates in time
and length units respectively, and \(ds^2\) has units of length squared.
Here \(c\) is the physical constant, not the Quadray component or PdF
edge length used in the volume sections.

\subsection{High-Precision Arithmetic
Note}\label{high-precision-arithmetic-note}

When evaluating determinants, FIMs, or geodesic distances for sensitive
problems, use quad precision (binary128) via GCC's
\passthrough{\lstinline!libquadmath!}
(\passthrough{\lstinline!\_\_float128!}, functions like
\passthrough{\lstinline!expq!}, \passthrough{\lstinline!sqrtq!}, and
\passthrough{\lstinline!quadmath\_snprintf!}). See
\href{https://gcc.gnu.org/onlinedocs/libquadmath/index.html}{GCC
libquadmath}. Where possible, it is useful to use symbolic math
libraries like SymPy to compute exact values.

\subsubsection{Reproducibility artifacts and external
validation}\label{reproducibility-artifacts-and-external-validation}

\begin{itemize}
\tightlist
\item
  \textbf{This manuscript's artifacts}: Raw data in
  \passthrough{\lstinline!quadmath/output/!} for reproducibility and
  downstream analysis:

  \begin{itemize}
  \tightlist
  \item
    \passthrough{\lstinline!fisher\_information\_matrix.csv!} /
    \passthrough{\lstinline!.npz!}: empirical Fisher matrix and inputs
  \item
    \passthrough{\lstinline!fisher\_information\_eigenvalues.csv!} /
    \passthrough{\lstinline!fisher\_information\_eigensystem.npz!}:
    eigenspectrum and eigenvectors
  \item
    \passthrough{\lstinline!natural\_gradient\_path.png!} with
    \passthrough{\lstinline!natural\_gradient\_path.csv!} /
    \passthrough{\lstinline!.npz!}: projected trajectory and raw
    coordinates
  \item
    \passthrough{\lstinline!ivm\_neighbors\_data.csv!} /
    \passthrough{\lstinline!ivm\_neighbors\_edges\_data.npz!}: neighbor
    coordinates (Quadray and XYZ)
  \item
    \passthrough{\lstinline!polyhedra\_quadray\_constructions.png!}:
    synergetics volume relationships schematic
  \end{itemize}
\item
  \textbf{External validation resources}: The
  \href{https://github.com/4dsolutions}{4dsolutions ecosystem} provides
  extensive cross-validation. See the \href{07_resources.md}{Resources}
  section for comprehensive details on computational implementations and
  validation.
\end{itemize}

\subsection{Namespaces summary
(notation)}\label{namespaces-summary-notation}

\begin{itemize}
\tightlist
\item
  Coxeter.4D: Euclidean E⁴; regular polytopes; not spacetime
  (cf.~Coxeter, Regular Polytopes, Dover ed., p.~119). Connections to
  higher-dimensional lattices and packings as in Conway \& Sloane.
\item
  Einstein.4D: Minkowski spacetime; indefinite metric; used here only as
  a metric analogy when discussing geodesics and information geometry.
\item
  Fuller.4D: Quadrays/IVM; tetrahedral lattice with integer tetravolume;
  unit regular tetrahedron has volume 1; synergetics scale relations
  (e.g., S3).
\end{itemize}

\end{document}
